% =============================================================================
\section{Problème étudié et modèle de trajectoire}
\label{sec:problem}
% =============================================================================

\subsection{Courbes de Bézier}

Une courbe de Bézier d'ordre $n$ est définie par $n+1$ points de contrôle
$\mathbf{P}_0, \ldots, \mathbf{P}_n$ et paramétrée par $t \in [0,1]$ :
\begin{equation}
\label{eq:bezier}
  \mathbf{B}(t) = \sum_{i=0}^{n} \binom{n}{i} (1-t)^{n-i}\, t^i \, \mathbf{P}_i.
\end{equation}
Les propriétés fondamentales sont :
\begin{itemize}[nosep]
  \item $\mathbf{B}(0) = \mathbf{P}_0$ et $\mathbf{B}(1) = \mathbf{P}_n$
        (interpolation des extrémités) ;
  \item la courbe est contenue dans l'enveloppe convexe des points de contrôle ;
  \item elle est infiniment dérivable, ce qui la rend adaptée à la modélisation
        de trajectoires fluides.
\end{itemize}

Dans notre problème, les points $\mathbf{P}_0$ (départ) et $\mathbf{P}_n$
(arrivée) sont fixés par l'instance.  Seuls les $n_{\text{CP}} = n - 1$ points
de contrôle internes constituent les \textbf{variables de décision}.

\subsection{Encodage des variables de décision}

L'encodage adopté est un vecteur plat $\mathbf{x} \in \mathbb{R}^d$ avec
$d = 2 \, n_{\text{CP}}$ :
\[
  \mathbf{x} = \bigl(x_0, y_0, \; x_1, y_1, \; \ldots, \; x_{n_{\text{CP}}-1},
                y_{n_{\text{CP}}-1}\bigr)
\]
où $(x_i, y_i)$ représente les coordonnées du $(i+1)$-ème point de contrôle
interne.  Les bornes de chaque composante sont celles du domaine rectangulaire
$[x_{\min}, x_{\max}] \times [y_{\min}, y_{\max}]$ défini par l'instance.

\subsection{Fonction de coût}
\label{sec:cost}

La fonction de coût $f(\mathbf{x})$ est évaluée en échantillonnant
$N = 1\,000$ points régulièrement espacés en paramètre $t \in [0,1]$ le long
de la courbe (nous notons $\mathbf{q}_k = \mathbf{B}(k/(N-1))$ pour
$k = 0, \ldots, N-1$), puis en sommant quatre composantes pondérées :

\begin{equation}
\label{eq:cost}
  f(\mathbf{x}) \;=\; L(\mathbf{x})
    \;+\; \alpha \cdot P_{\text{obs}}(\mathbf{x})
    \;+\; \beta \cdot P_{\text{courb}}(\mathbf{x})
    \;+\; \gamma \cdot P_{\text{bord}}(\mathbf{x})
\end{equation}
avec $\alpha = 100$, $\beta = 1$, $\gamma = 100$.

\paragraph{Longueur $L$.}
Somme des distances euclidiennes entre points consécutifs :
\[
  L(\mathbf{x}) = \sum_{k=1}^{N-1} \|\mathbf{q}_k - \mathbf{q}_{k-1}\|_2.
\]

\paragraph{Pénalité d'obstacle $P_{\text{obs}}$.}
Pour chaque obstacle circulaire de centre $\mathbf{c}_j$ et rayon $r_j$, et
pour chaque point de trajectoire $\mathbf{q}_k$, la distance
$d_{jk} = \|\mathbf{q}_k - \mathbf{c}_j\|_2$ détermine la pénalité :
\[
  p_{\text{obs}}(d, r) =
  \begin{cases}
    2 + r - d, & \text{si } d \leq r \quad \text{(collision dure)}, \\
    1 - (d - r), & \text{si } r < d < r+1 \quad \text{(zone de proximité)}, \\
    0, & \text{sinon}.
  \end{cases}
\]
La pénalité totale est $P_{\text{obs}} = \sum_j \sum_k p_{\text{obs}}(d_{jk}, r_j)$.

\emph{Discussion.}  Cette formulation introduit une \textbf{zone de transition
douce} (anneau de largeur 1 autour de chaque obstacle) qui lisse le paysage
d'optimisation.  Sans cette zone, la frontière abrupte entre « collision » et
« pas de collision » créerait des discontinuités très défavorables aux
métaheuristiques à population.

\paragraph{Pénalité de courbure $P_{\text{courb}}$.}
Somme des déviations angulaires entre triplets de points consécutifs :
\[
  P_{\text{courb}} = \sum_{k=1}^{N-2} \bigl(1 - \cos\theta_k\bigr),
  \qquad
  \cos\theta_k = \frac{(\mathbf{q}_k - \mathbf{q}_{k-1}) \cdot
                        (\mathbf{q}_{k+1} - \mathbf{q}_k)}
                       {\|\mathbf{q}_k - \mathbf{q}_{k-1}\|\;
                        \|\mathbf{q}_{k+1} - \mathbf{q}_k\|}.
\]

\paragraph{Pénalité de sortie $P_{\text{bord}}$.}
Nombre de points hors du domaine :
$P_{\text{bord}} = \#\{k : \mathbf{q}_k \notin
[x_{\min}, x_{\max}] \times [y_{\min}, y_{\max}]\}$.

\emph{Discussion.}  Les pondérations $\alpha = \gamma = 100$ et $\beta = 1$
sont imposées par le cadre du projet (API \texttt{Problem} fournie par
l'enseignant).  Elles rendent les collisions et les sorties de zone cent fois
plus coûteuses par unité que la longueur, assurant que la faisabilité
(absence de collision) prime sur l'optimalité (longueur minimale).  La courbure
($\beta = 1$) agit comme un régularisateur léger favorisant la fluidité sans
dominer les autres termes.  Ce choix de hiérarchie est cohérent avec les
recommandations de la littérature sur les méthodes de pénalité
extérieure~\cite{coello2002}.

\subsection{Description des instances}
\label{sec:instances}

Les quatre instances fournies sont résumées dans le tableau~\ref{tab:instances}.

\begin{table}[H]
\centering
\caption{Caractéristiques des quatre instances du problème.}
\label{tab:instances}
\begin{tabular}{@{}lccccl@{}}
\toprule
Instance & $n_{\text{CP}}$ & $d$ & Obstacles & Départ $\to$ Arrivée & Description \\
\midrule
\texttt{prob1} & 16 & 32 & 5  & $(1,1) \to (25,25)$  & Diagonale, obstacles épars \\
\texttt{prob2} & 4  & 8  & 1  & $(11,1) \to (11,21)$  & Vertical, gros obstacle ($r=9$) \\
\texttt{prob3} & 8  & 16 & 33 & $(1,13) \to (25,13)$ & 3 murs verticaux, brèches \\
\texttt{prob4} & 8  & 16 & 33 & $(1,1) \to (25,1)$   & Labyrinthe en S, 3 murs \\
\bottomrule
\end{tabular}
\end{table}

\emph{Discussion.}  La difficulté croît de \texttt{prob1} (peu d'obstacles,
haute dimension mais paysage lisse) à \texttt{prob4} (33~obstacles formant
trois murs avec des brèches alternées haut/bas, nécessitant un chemin en S
que la courbe de Bézier de degré~10 peine à reproduire).
